\documentclass[11pt]{article}

\usepackage[margin=1in]{geometry}
\usepackage[T1]{fontenc}
\usepackage[utf8]{inputenc}
\usepackage{hyperref}
\usepackage{xurl}
\usepackage{enumitem}
\usepackage{longtable}
\usepackage{xcolor}
\usepackage{array}

\hypersetup{
  colorlinks=true,
  linkcolor=blue,
  urlcolor=blue,
  citecolor=blue
}

\setlist[itemize]{leftmargin=1.5em}
\setlist[enumerate]{leftmargin=1.7em}
\renewcommand{\arraystretch}{1.25}

\title{Abshaar: Step-by-Step Implementation Guide}
\author{Rauf Nawaz}
\date{May 23, 2026}

\begin{document}
\maketitle
\tableofcontents
\newpage

\section{Purpose of This Guide}

This guide explains how to build Abshaar from an empty repository into a working
open-source system for translating and explaining Punjabi, Urdu, Persian-
influenced, Braj, Bhakti, and Sufi poetry in English. It is intentionally
granular. The goal is to make the project understandable for future contributors
who may care deeply about poetry but may not be machine-learning engineers.

The project should be built in this order:

\begin{enumerate}
  \item Build the corpus and annotation structure.
  \item Create a small gold-standard dataset by hand.
  \item Build a local AI-assisted draft pipeline.
  \item Add retrieval over glossary, context, and prior reviewed examples.
  \item Build a static public website.
  \item Add a review loop.
  \item Fine-tune only after enough reviewed examples exist.
  \item Add chatbot features only when the archive itself is strong.
\end{enumerate}

\subsection{Core Principle}

Do not begin by training a model. Begin by building a careful archive. The model
will only become useful if the project contains strong source text, source
metadata, translations, glossary entries, review notes, and poet-specific
context.

\section{What You Are Building}

Abshaar should have three connected layers.

\subsection{Layer 1: Archive}

The archive stores the original poem, transliteration, source information,
translation, explanation, glossary, poet context, timeline events, and review
status.

\subsection{Layer 2: Local AI Assistant}

The AI assistant drafts literal translations, literary translations, tashreeh,
glossary notes, and question-answer responses. It must retrieve relevant context
before answering. It must never replace human review.

\subsection{Layer 3: Public Website}

The website lets readers explore poets, poems, timelines, themes, glossary terms,
and beginner-friendly explanations. It should run for free on GitHub Pages in the
first public version.

\section{High-Level Timeline}

\begin{longtable}{>{\raggedright\arraybackslash}p{0.16\textwidth}
                  >{\raggedright\arraybackslash}p{0.28\textwidth}
                  >{\raggedright\arraybackslash}p{0.45\textwidth}}
\textbf{Phase} & \textbf{Time} & \textbf{Outcome} \\
\hline
0 & 1--2 weeks & Choose one poet, source rules, licensing, and annotation standards. \\
1 & 2--4 weeks & Build a small corpus of 20--50 works with source metadata. \\
2 & 3--5 weeks & Build local draft pipeline using existing open models. \\
3 & Ongoing & Store human corrections as structured review data. \\
4 & 2 weeks & Create evaluation set and scoring rubric. \\
5 & After 500+ reviews & Try LoRA or QLoRA fine-tuning. \\
6 & 3--6 weeks & Build static website MVP. \\
7 & After website MVP & Add contribution workflow and public review process. \\
8 & Long term & Expand poets, themes, dictionary, audio, and optional chatbot. \\
\end{longtable}

\section{Phase 0: Project Decisions}

\subsection{Step 0.1: Choose the First Poet}

Choose one poet for the first version. Do not start with all Sufi and Bhakti
poetry at once.

Recommended first poet options:

\begin{itemize}
  \item Kabir: rich interpretive tradition, many public-domain materials, strong
  cross-religious vocabulary.
  \item Bulleh Shah: central for Punjabi Sufi thought, strong theme of ishq,
  murshid, social critique, and performative tradition.
  \item Baba Farid: important for Punjabi, Sikh scripture context, and Sufi
  vocabulary.
  \item Shah Hussain: strong for Punjabi kafi form and mystical love.
  \item Waris Shah: larger literary project, better after the system matures.
\end{itemize}

Recommended first choice for Abshaar: Bulleh Shah. He fits the project's Punjabi,
Urdu, Persianate, and Sufi focus especially well. His poetry also forces the
system to handle the core problem: literal translation is not enough because
words such as ishq, murshid, nafs, haq, and kafir can shift meaning depending on
poetic, spiritual, and social context.

Begin with 20--50 short kafis or selected stanzas. Do not begin with a large
complete collected works edition. The first release should show quality, not
scale.

\subsection{Step 0.2: Decide the First Release Scope}

Write down the first release scope in a file called \verb|docs/mvp_scope.md|.

Include:

\begin{itemize}
  \item selected poet;
  \item number of works;
  \item allowed source texts;
  \item scripts to preserve;
  \item transliteration standard;
  \item translation style;
  \item explanation style;
  \item what is explicitly out of scope.
\end{itemize}

Example:

\begin{verbatim}
MVP poet: Bulleh Shah
Works: 30 short couplets or songs
Scripts: Shahmukhi/Perso-Arabic where available, Gurmukhi or other source
scripts if using a licensed source, plus Latin transliteration
Translation types: literal gloss, literary translation, tashreeh
Website: static pages only, no public chatbot yet
Model: local Qwen3 through Ollama for draft explanations
Review: all publishable outputs require human approval
\end{verbatim}

\subsection{Step 0.3: Decide Licensing}

Use separate licenses for code and content.

\begin{itemize}
  \item Code: MIT or Apache-2.0.
  \item Project documentation: CC BY 4.0.
  \item Your translations, tashreeh, glossary, and annotations: CC BY-SA 4.0.
  \item Source texts: keep their own license and source status.
\end{itemize}

Important: do not assume that modern translations, recordings, transcriptions,
or website annotations are free to reuse. Classical poetry may be public domain,
but a modern edition or translation may still be copyrighted.

\subsection{Step 0.4: Define the Translation Philosophy}

Create a short translation style guide. The model and contributors need to know
what ``good'' means.

Recommended rules:

\begin{enumerate}
  \item Preserve the image before explaining the abstraction.
  \item Keep literal translation separate from literary translation.
  \item Keep tashreeh separate from translation.
  \item Preserve culturally loaded terms when English would flatten them.
  \item Use footnotes or glossary entries instead of forcing everything into the
  English line.
  \item Mark uncertainty and alternate readings.
  \item Do not collapse Sufi, Bhakti, Sikh, Islamic, Hindu, folk, and regional
  concepts into one generic mystical vocabulary.
\end{enumerate}

\section{Phase 1: Repository and Local Setup}

\subsection{Step 1.1: Confirm Git Is Working}

From PowerShell:

\begin{verbatim}
cd "D:\Harvard\Poetry Model Project"
git status --short --branch
\end{verbatim}

If the folder is not a Git repository, initialize it:

\begin{verbatim}
git init
git branch -M main
\end{verbatim}

If the remote is not set:

\begin{verbatim}
git remote add origin https://github.com/RaufNawaz/Abshaar.git
\end{verbatim}

\subsection{Step 1.2: Keep a Clean Folder Structure}

Use this structure:

\begin{verbatim}
docs/
  01_project_roadmap.md
  02_model_strategy.md
  03_data_and_annotation_guide.md
  04_website_architecture.md
  05_local_setup.md
  06_open_source_governance.md
  07_step_by_step_implementation_guide.tex

data/
  raw/
    public/
    private/
  processed/
  annotations/
  lexicon/
  context/
  cache/

scripts/
  ingest/
  translate/
  review/
  export/

src/
  abshaar/
    __init__.py
    schemas.py
    ingest.py
    translate.py
    retrieve.py
    review.py
    export_site.py

website/
\end{verbatim}

Do not put model files, private scans, or generated caches into Git.

\subsection{Step 1.3: Install Required Software}

Install these:

\begin{itemize}
  \item Git;
  \item Python 3.11 or 3.12;
  \item Node.js LTS;
  \item VS Code;
  \item Ollama;
  \item optional: WSL2 with Ubuntu for heavier ML work.
\end{itemize}

\subsection{Step 1.4: Create Python Environment}

\begin{verbatim}
cd "D:\Harvard\Poetry Model Project"
py -0p
py -m venv .venv
.\.venv\Scripts\Activate.ps1
python -m pip install --upgrade pip
\end{verbatim}

Use Python 3.11 or newer. The command \verb|py -0p| shows which Python versions
are installed on the machine.

Install minimal dependencies first:

\begin{verbatim}
pip install pydantic pandas rich
\end{verbatim}

Add ML dependencies when needed:

\begin{verbatim}
pip install torch transformers sentence-transformers accelerate
pip install chromadb fastapi uvicorn
\end{verbatim}

Add fine-tuning tools later:

\begin{verbatim}
pip install peft datasets evaluate
\end{verbatim}

\section{Phase 2: Build the Corpus MVP}

\subsection{Step 2.1: Choose Safe Source Texts}

For each source, record:

\begin{itemize}
  \item title;
  \item editor or translator if applicable;
  \item publication date;
  \item URL or bibliographic reference;
  \item copyright status;
  \item whether it may be used for training;
  \item whether it may be republished on the website.
\end{itemize}

Create \verb|data/context/sources.jsonl|. Each line should look like:

\begin{verbatim}
{
  "id": "source_0001",
  "title": "Source title",
  "author_or_editor": "Name",
  "year": "1910",
  "url": "https://example.org",
  "rights_status": "public-domain",
  "allowed_uses": ["quote", "train", "republish"],
  "notes": "Why this is safe to use."
}
\end{verbatim}

\subsection{Step 2.2: Create the Poem Schema}

Create \verb|data/processed/poems.jsonl|. Each line is one poem record:

\begin{verbatim}
{
  "id": "kabir_0001",
  "poet_id": "kabir",
  "title": "First line or working title",
  "work_type": "couplet",
  "source_ids": ["source_0001"],
  "original": {
    "text": "Original text here",
    "script": "Devanagari",
    "language_spans": [
      {
        "text": "Original span",
        "language": "Braj",
        "confidence": 0.8
      }
    ]
  },
  "transliteration": {
    "scheme": "project-latin-v1",
    "text": "Transliteration here"
  },
  "translations": [],
  "tashreeh": [],
  "glossary_terms": [],
  "themes": [],
  "review_status": "draft"
}
\end{verbatim}

\subsection{Step 2.3: Start With 20 Works}

For each work:

\begin{enumerate}
  \item Add original text.
  \item Add source ID.
  \item Add transliteration.
  \item Add a literal gloss.
  \item Add your first literary translation.
  \item Add tashreeh.
  \item Add glossary terms.
  \item Add themes.
  \item Mark review status.
\end{enumerate}

Do not let the model create the gold dataset alone. The first 20 works should be
human-led. This gives the system a standard to imitate later.

\subsection{Step 2.4: Build the Lexicon}

Create \verb|data/lexicon/terms.jsonl|. Every term should be poet-specific when
needed.

\begin{verbatim}
{
  "id": "term_ishq_bulleh_shah",
  "headword": "ishq",
  "transliteration": "ishq",
  "languages": ["Persian", "Urdu", "Punjabi"],
  "basic_meaning": "love",
  "poet_specific_meaning": "Divine longing that dissolves ego and rank.",
  "do_not_flatten_to": ["romantic love"],
  "related_terms": ["fana", "murshid", "haqiqat"],
  "example_poems": ["bulleh_0001"],
  "source_ids": ["source_0001"],
  "review_status": "draft"
}
\end{verbatim}

Add a field for terms that should remain untranslated. Example: ``naam'',
``shabad'', ``murshid'', ``haq'', ``fana'', ``ishq'', ``virah''. Sometimes the
right English move is to preserve the original word and explain it.

\subsection{Step 2.5: Build Poet Context}

Create:

\begin{itemize}
  \item \verb|data/context/people.jsonl|;
  \item \verb|data/context/events.jsonl|;
  \item \verb|data/context/themes.jsonl|.
\end{itemize}

For timeline events:

\begin{verbatim}
{
  "id": "event_kabir_0001",
  "poet_id": "kabir",
  "date_label": "15th century",
  "date_start": null,
  "date_end": null,
  "event_type": "historical_context",
  "title": "North Indian devotional and Sufi milieu",
  "description": "Short, sourced note.",
  "source_ids": ["source_0002"],
  "confidence": 0.6
}
\end{verbatim}

Use confidence because many classical poet biographies are uncertain.

\section{Phase 3: Create the Human Review Workflow}

\subsection{Step 3.1: Define Review Status Values}

Use these statuses:

\begin{itemize}
  \item \verb|draft|;
  \item \verb|machine_draft|;
  \item \verb|human_reviewed|;
  \item \verb|needs_source_check|;
  \item \verb|needs_language_check|;
  \item \verb|publishable|;
  \item \verb|rejected|.
\end{itemize}

\subsection{Step 3.2: Create Review Records}

Create \verb|data/annotations/reviews.jsonl|.

\begin{verbatim}
{
  "id": "review_0001",
  "poem_id": "kabir_0001",
  "reviewer": "Rauf",
  "date": "2026-05-23",
  "scores": {
    "source_fidelity": 4,
    "metaphor_fidelity": 3,
    "poet_specific_context": 3,
    "literary_quality": 4,
    "beginner_clarity": 5,
    "humility_about_uncertainty": 4
  },
  "problems": [
    {
      "type": "missed_context",
      "note": "The image was translated literally but not explained."
    }
  ],
  "corrected_translation": "Corrected translation here.",
  "corrected_tashreeh": "Corrected explanation here.",
  "style_note": "Preserve the metaphor before explaining it."
}
\end{verbatim}

\subsection{Step 3.3: Make Review a Habit}

After every AI-generated output, ask:

\begin{enumerate}
  \item Did it preserve the concrete image?
  \item Did it understand poet-specific vocabulary?
  \item Did it invent context?
  \item Did it flatten religious differences?
  \item Did it make the English readable?
  \item Did it explain too much?
  \item Did it explain too little?
  \item Is it safe to publish?
\end{enumerate}

The review data is the project's most valuable machine-learning asset.

\section{Phase 4: Build the First Local AI Pipeline}

\subsection{Step 4.1: Model Roles}

Use different models for different roles.

\begin{longtable}{>{\raggedright\arraybackslash}p{0.22\textwidth}
                  >{\raggedright\arraybackslash}p{0.30\textwidth}
                  >{\raggedright\arraybackslash}p{0.38\textwidth}}
\textbf{Task} & \textbf{Recommended Model} & \textbf{Reason} \\
\hline
Baseline translation & IndicTrans2 & Strong Indic-to-English baseline for supported languages. \\
Fallback translation & NLLB-200 & Broader language coverage, useful as fallback. \\
Explanation and rewriting & Qwen3 via Ollama & Strong local instruction-following and multilingual ability. \\
Retrieval embeddings & BGE-M3 & Multilingual embeddings for glossary and context retrieval. \\
Website search & Pagefind or MiniSearch & Static, free, no backend required. \\
\end{longtable}

\subsection{Step 4.2: Install and Test Ollama}

Ollama is the simplest local model runner for the first version. It gives you a
local HTTP API, which means the rest of the project can call the model without
special cloud credentials.

Install Ollama from:

\begin{verbatim}
https://ollama.com/
\end{verbatim}

After installation, confirm it is running:

\begin{verbatim}
ollama --version
ollama list
\end{verbatim}

If the Ollama background service is running, this URL should respond:

\begin{verbatim}
http://localhost:11434
\end{verbatim}

\subsection{Step 4.3: Install the Interpretive LLM}

\begin{verbatim}
ollama pull qwen3:8b
ollama run qwen3:8b
\end{verbatim}

If the machine is slow:

\begin{verbatim}
ollama pull qwen3:4b
\end{verbatim}

If the machine is stronger:

\begin{verbatim}
ollama pull qwen3:14b
\end{verbatim}

Use Qwen3 for:

\begin{itemize}
  \item tashreeh drafting;
  \item literary translation rewriting;
  \item glossary explanation;
  \item alternate readings;
  \item question-answering over retrieved context;
  \item self-checking whether an explanation overclaims.
\end{itemize}

Do not use Qwen3 alone as the translation engine. Use it after baseline
translation and retrieval.

\subsection{Step 4.4: Call Ollama From Python}

Install the Python client:

\begin{verbatim}
pip install ollama
\end{verbatim}

Minimal test:

\begin{verbatim}
from ollama import chat

response = chat(
    model="qwen3:8b",
    messages=[
        {"role": "user", "content": "Explain what tashreeh means in one sentence."}
    ],
)

print(response["message"]["content"])
\end{verbatim}

For project use, wrap this in a function:

\begin{verbatim}
def ask_llm(system_prompt, user_prompt, model="qwen3:8b"):
    response = chat(
        model=model,
        messages=[
            {"role": "system", "content": system_prompt},
            {"role": "user", "content": user_prompt},
        ],
        options={
            "temperature": 0.3,
            "top_p": 0.9,
        },
    )
    return response["message"]["content"]
\end{verbatim}

\subsection{Step 4.5: Install Baseline Translation Models}

Use two baseline translation routes.

\subsubsection{IndicTrans2}

Use IndicTrans2 for supported Indic-to-English translation. It is the preferred
baseline for supported Indic languages, including Punjabi and Urdu-related
coverage where available.

Install dependencies:

\begin{verbatim}
pip install torch transformers IndicTransToolkit
\end{verbatim}

Model to start with:

\begin{verbatim}
ai4bharat/indictrans2-indic-en-1B
\end{verbatim}

Important:

\begin{itemize}
  \item some Hugging Face model files may require accepting terms on the model
  page;
  \item do not commit model weights to Git;
  \item use this model for baseline meaning, not final literary translation;
  \item for poetry, line segmentation matters.
\end{itemize}

\subsubsection{NLLB-200}

Use NLLB-200 as a fallback when IndicTrans2 does not handle the source well.

Install:

\begin{verbatim}
pip install sentencepiece sacremoses
\end{verbatim}

Model:

\begin{verbatim}
facebook/nllb-200-distilled-600M
\end{verbatim}

Use NLLB for:

\begin{itemize}
  \item fallback baseline translation;
  \item comparison against IndicTrans2;
  \item broader multilingual snippets;
  \item debugging whether a word/phrase is being lost.
\end{itemize}

\subsection{Step 4.6: Install Embedding and Retrieval Models}

Use BGE-M3 for multilingual retrieval.

\begin{verbatim}
pip install sentence-transformers chromadb
\end{verbatim}

Minimal embedding test:

\begin{verbatim}
from sentence_transformers import SentenceTransformer

model = SentenceTransformer("BAAI/bge-m3")
texts = [
    "Ishq in Bulleh Shah often points to divine love and ego-dissolution.",
    "Murshid refers to the spiritual guide in Sufi contexts."
]
embeddings = model.encode(texts)
print(embeddings.shape)
\end{verbatim}

For the first version, store embeddings in Chroma as a generated cache. Keep the
canonical source data in JSONL.

\subsection{Step 4.7: How the Models Work Together}

The project should use the models in a pipeline, not as isolated chatbots.

\begin{enumerate}
  \item Load one Bulleh Shah poem record.
  \item Detect script and segment the poem into lines.
  \item Run IndicTrans2 for baseline translation if the language/script is
  supported.
  \item Run NLLB-200 as a fallback or comparison.
  \item Retrieve glossary entries using exact match and BGE-M3 semantic search.
  \item Retrieve poet context, themes, and reviewed examples.
  \item Send the original, transliteration, baseline translations, and retrieved
  context to Qwen3.
  \item Ask Qwen3 for structured output: literal gloss, literary translation,
  tashreeh, key terms, alternate readings, and uncertainty.
  \item Save the raw output.
  \item Review manually.
  \item Store the corrected version as the new gold example.
\end{enumerate}

\subsection{Step 4.8: Recommended Python Modules}

Create these files later:

\begin{verbatim}
src/abshaar/schemas.py
src/abshaar/io.py
src/abshaar/translate_baseline.py
src/abshaar/embed.py
src/abshaar/retrieve.py
src/abshaar/generate.py
src/abshaar/review.py
src/abshaar/export_site.py
\end{verbatim}

Responsibilities:

\begin{itemize}
  \item \verb|schemas.py|: Pydantic models for poems, terms, sources, reviews.
  \item \verb|io.py|: read and write JSONL.
  \item \verb|translate_baseline.py|: IndicTrans2 and NLLB wrappers.
  \item \verb|embed.py|: BGE-M3 embedding generation.
  \item \verb|retrieve.py|: exact match plus vector retrieval.
  \item \verb|generate.py|: Qwen3 prompts through Ollama.
  \item \verb|review.py|: save human corrections.
  \item \verb|export_site.py|: export clean public data for the website.
\end{itemize}

\subsection{Step 4.9: Draft Pipeline Steps}

The first pipeline should run like this:

\begin{enumerate}
  \item Load one poem from \verb|poems.jsonl|.
  \item Segment the poem into lines or couplets.
  \item Detect script and probable language spans.
  \item Create or load transliteration.
  \item Run baseline translation.
  \item Retrieve glossary entries and poet context.
  \item Ask the LLM to produce a structured draft.
  \item Save the model output.
  \item Ask a human reviewer to correct it.
  \item Save the review.
\end{enumerate}

\subsection{Step 4.10: Required Output Format}

The model should return JSON-like structured data:

\begin{verbatim}
{
  "literal_gloss": "...",
  "direct_translation": "...",
  "literary_translation": "...",
  "tashreeh_beginner": "...",
  "tashreeh_advanced": "...",
  "key_terms": [
    {
      "term": "...",
      "meaning_in_context": "...",
      "uncertainty": "low/medium/high"
    }
  ],
  "alternate_readings": ["..."],
  "uncertainty_notes": ["..."],
  "retrieved_context_used": ["context_id_1", "term_id_2"]
}
\end{verbatim}

If a model cannot return valid JSON reliably, store the raw output and parse it
later. Do not lose raw outputs; they are useful for debugging.

\subsection{Step 4.11: Prompt Template}

Use a staged prompt, not a vague request.

\begin{verbatim}
You are assisting with a scholarly but beginner-friendly translation
of classical South Asian mystical poetry.

Rules:
1. Preserve the original image before explaining it.
2. Separate literal gloss, literary translation, and tashreeh.
3. Do not invent biography or doctrine.
4. Use retrieved context only when relevant.
5. Mark uncertainty and alternate readings.
6. If a term should remain untranslated, say so.

Poet:
{poet_name}

Original:
{original_text}

Transliteration:
{transliteration}

Known glossary:
{retrieved_glossary}

Poet context:
{retrieved_context}

Prior reviewed examples:
{few_shot_examples}

Return:
- literal_gloss
- direct_translation
- literary_translation
- beginner_tashreeh
- advanced_tashreeh
- key_terms
- alternate_readings
- uncertainty_notes
\end{verbatim}

\section{Phase 4A: Initial Dataset Template for Bulleh Shah}

\subsection{Why Bulleh Shah First}

Bulleh Shah is the recommended first poet because he is directly aligned with the
project's goals:

\begin{itemize}
  \item Punjabi/Sufi center of gravity;
  \item Persian and Islamic vocabulary woven into Punjabi poetic expression;
  \item strong metaphysical concepts that cannot be translated literally;
  \item social critique, religious critique, love, ego, and spiritual authority;
  \item strong usefulness for a future public glossary.
\end{itemize}

\subsection{Initial Dataset Files}

Create template files under:

\begin{verbatim}
data/templates/
\end{verbatim}

Then copy them into the real data folders when you begin:

\begin{verbatim}
data/templates/poems.template.jsonl        -> data/processed/poems.jsonl
data/templates/terms.template.jsonl        -> data/lexicon/terms.jsonl
data/templates/sources.template.jsonl      -> data/context/sources.jsonl
data/templates/people.template.jsonl       -> data/context/people.jsonl
data/templates/events.template.jsonl       -> data/context/events.jsonl
data/templates/themes.template.jsonl       -> data/context/themes.jsonl
data/templates/reviews.template.jsonl      -> data/annotations/reviews.jsonl
data/templates/model_outputs.template.jsonl -> data/annotations/model_outputs.jsonl
\end{verbatim}

\subsection{Minimum Fields for Every Bulleh Shah Poem}

Every poem record should include:

\begin{itemize}
  \item stable ID, such as \verb|bulleh_shah_0001|;
  \item source IDs;
  \item title or first line;
  \item original text;
  \item script;
  \item language spans;
  \item transliteration;
  \item segmentation;
  \item literal gloss;
  \item literary translation;
  \item tashreeh;
  \item key terms;
  \item themes;
  \item review status;
  \item rights status.
\end{itemize}

\subsection{Starter Glossary Terms for Bulleh Shah}

Begin with terms like:

\begin{itemize}
  \item ishq;
  \item murshid;
  \item pir;
  \item haq;
  \item nafs;
  \item fana;
  \item kafir;
  \item masjid;
  \item mandir;
  \item yaar;
  \item rang;
  \item faqir.
\end{itemize}

Each term should have a basic meaning and a Bulleh Shah-specific meaning. The
poet-specific meaning is what will make the model useful.

\section{Phase 5: Retrieval-Augmented Generation}

\subsection{Step 5.1: What Should Be Retrieved}

For each poem, retrieve:

\begin{itemize}
  \item glossary terms that appear in the poem;
  \item poet-specific notes;
  \item related reviewed poems;
  \item theme notes;
  \item timeline events relevant to the poem;
  \item source notes.
\end{itemize}

\subsection{Step 5.2: Create Embedding Documents}

Transform records into small searchable passages.

Example glossary embedding document:

\begin{verbatim}
ID: term_ishq_bulleh_shah
TYPE: glossary
TEXT:
For Bulleh Shah, ishq often means divine longing and transformative love
that challenges social hierarchy, ego, and formal identity. It should not
be reduced to romantic love without context.
\end{verbatim}

Example timeline embedding document:

\begin{verbatim}
ID: event_kabir_0001
TYPE: event
TEXT:
Kabir is associated with a North Indian devotional world in which Hindu,
Islamic, artisan, and oral traditions intersect. Biographical certainty is
limited, so claims should be presented carefully.
\end{verbatim}

\subsection{Step 5.3: Retrieval Rules}

For every model answer:

\begin{enumerate}
  \item retrieve at least three relevant context chunks;
  \item pass them into the prompt;
  \item ask the model to cite which chunks it used;
  \item reject answers that make claims not grounded in the source or marked as
  interpretation;
  \item store retrieval IDs with the generated output.
\end{enumerate}

\section{Phase 6: Evaluation Before Fine-Tuning}

\subsection{Step 6.1: Build an Evaluation Set}

Create 50--100 passages that are held out from normal prompting. These should not
be repeatedly used as few-shot examples.

Each evaluation item should include:

\begin{itemize}
  \item original text;
  \item transliteration;
  \item human literal gloss;
  \item human literary translation;
  \item human tashreeh;
  \item key terms;
  \item accepted alternate readings;
  \item source IDs.
\end{itemize}

\subsection{Step 6.2: Scoring Rubric}

Score every model output from 1 to 5:

\begin{longtable}{>{\raggedright\arraybackslash}p{0.26\textwidth}
                  >{\raggedright\arraybackslash}p{0.64\textwidth}}
\textbf{Category} & \textbf{Question} \\
\hline
Source fidelity & Does it preserve what the original says? \\
Image fidelity & Does it preserve the metaphor before abstracting it? \\
Poet-specific context & Does it use this poet's meaning of key terms? \\
Cultural grounding & Does it explain context without overclaiming? \\
Literary quality & Does the English feel alive, not mechanical? \\
Beginner clarity & Can a new reader understand the stakes? \\
Uncertainty & Does it mark ambiguity and avoid false certainty? \\
Citation discipline & Does it stay grounded in source and retrieved context? \\
\end{longtable}

\subsection{Step 6.3: Decide If Fine-Tuning Is Worth It}

Do not fine-tune until:

\begin{itemize}
  \item you have at least 500 reviewed examples;
  \item the prompt-and-retrieval system is already working;
  \item you know exactly what behavior needs improvement;
  \item you have a held-out evaluation set;
  \item you can compare pre-fine-tune and post-fine-tune outputs.
\end{itemize}

\section{Phase 7: Fine-Tuning Plan}

\subsection{Step 7.1: Fine-Tune the Right Thing}

Fine-tuning should teach:

\begin{itemize}
  \item output format;
  \item translation style;
  \item explanation style;
  \item how to preserve terms;
  \item how to mark uncertainty;
  \item how to avoid flattening metaphors.
\end{itemize}

Fine-tuning should not be the main storage place for knowledge. Knowledge should
stay in the corpus, glossary, and retrieval system.

\subsection{Step 7.2: Training Data Format}

Use instruction-response records:

\begin{verbatim}
{
  "messages": [
    {
      "role": "system",
      "content": "You translate and explain classical mystical poetry..."
    },
    {
      "role": "user",
      "content": "Poet: Kabir\nOriginal: ...\nContext: ..."
    },
    {
      "role": "assistant",
      "content": "{ reviewed JSON output here }"
    }
  ]
}
\end{verbatim}

\subsection{Step 7.3: LoRA or QLoRA}

Use LoRA or QLoRA rather than full fine-tuning.

Recommended progression:

\begin{enumerate}
  \item Try prompt-only baseline.
  \item Try prompt plus retrieval.
  \item Try few-shot examples.
  \item Try LoRA on 500--2,000 examples.
  \item Try preference tuning only after collecting rejected and corrected pairs.
\end{enumerate}

\subsection{Step 7.4: Training Safety}

Before releasing a fine-tuned model:

\begin{itemize}
  \item confirm training data rights;
  \item remove private notes;
  \item remove copyrighted modern translations unless licensed;
  \item include a model card;
  \item disclose base model;
  \item disclose intended use and limitations;
  \item explain that outputs require human review.
\end{itemize}

\section{Phase 8: Website Development}

\subsection{Step 8.1: Start Static}

Use Astro for the first website.

\begin{verbatim}
npm create astro@latest website
cd website
npm install
npm run dev
\end{verbatim}

Use a static site because:

\begin{itemize}
  \item it can be hosted on GitHub Pages for free;
  \item it is safer and easier to maintain;
  \item contributors can edit content without managing servers;
  \item it keeps the archive public even if no AI backend is running.
\end{itemize}

\subsection{Step 8.2: Website Pages}

Build pages in this order:

\begin{enumerate}
  \item Home page with search and featured poet.
  \item Poet profile page.
  \item Works index.
  \item Poem page.
  \item Glossary page.
  \item Theme page.
  \item Timeline page.
  \item Source and credits page.
  \item Contribute page.
  \item Ask page with curated FAQ.
\end{enumerate}

\subsection{Step 8.3: Poem Page Requirements}

Each poem page should show:

\begin{itemize}
  \item original text;
  \item transliteration;
  \item literal gloss;
  \item literary translation;
  \item beginner tashreeh;
  \item advanced notes if available;
  \item glossary highlights;
  \item themes;
  \item related works;
  \item source and license;
  \item review status.
\end{itemize}

\subsection{Step 8.4: Timeline Requirements}

Timeline entries should distinguish:

\begin{itemize}
  \item life event;
  \item uncertain biography;
  \item historical context;
  \item composition or oral tradition;
  \item manuscript or publication;
  \item later reception;
  \item modern performance or interpretation.
\end{itemize}

Use confidence labels. Classical poet timelines are often uncertain, and the site
should not pretend otherwise.

\subsection{Step 8.5: Chatbot Options}

Do not make the chatbot the first public feature. Build search and curated FAQ
first.

Options:

\begin{enumerate}
  \item Static FAQ and search: best first release.
  \item Browser-local AI through WebLLM: free but hardware-dependent.
  \item Separate FastAPI backend with Ollama: best quality but requires hosting.
\end{enumerate}

GitHub Pages alone cannot host a normal server-side chatbot.

\section{Phase 9: Public Contribution Workflow}

\subsection{Step 9.1: Contributor Paths}

Allow contributors to help without understanding the whole system.

Contribution types:

\begin{itemize}
  \item source correction;
  \item transliteration correction;
  \item glossary addition;
  \item alternate translation;
  \item tashreeh improvement;
  \item timeline event;
  \item theme note;
  \item website bug fix.
\end{itemize}

\subsection{Step 9.2: Contribution Template}

Every content contribution should answer:

\begin{enumerate}
  \item What are you changing?
  \item What source supports it?
  \item What is the rights status?
  \item Is there uncertainty?
  \item Does this affect published pages?
  \item Does this require language review?
  \item Does this require interpretation review?
\end{enumerate}

\subsection{Step 9.3: Review Roles}

Use these roles:

\begin{itemize}
  \item maintainer;
  \item language reviewer;
  \item interpretation reviewer;
  \item source/license reviewer;
  \item technical reviewer.
\end{itemize}

At the beginning, one person may hold several roles. Over time, separate them.

\section{Detailed Weekly Plan}

\subsection{Week 1}

\begin{itemize}
  \item Finalize MVP poet.
  \item Create \verb|docs/mvp_scope.md|.
  \item Create source rules.
  \item Create translation style guide.
  \item Confirm licensing.
  \item Confirm GitHub remote and branch workflow.
\end{itemize}

\subsection{Week 2}

\begin{itemize}
  \item Collect 20 safe source texts.
  \item Create \verb|sources.jsonl|.
  \item Create first \verb|poems.jsonl| records.
  \item Add transliterations.
  \item Add source notes.
\end{itemize}

\subsection{Week 3}

\begin{itemize}
  \item Write literal glosses for 10 poems.
  \item Write literary translations for 10 poems.
  \item Write beginner tashreeh for 10 poems.
  \item Add 20 glossary terms.
  \item Add 10 theme records.
\end{itemize}

\subsection{Week 4}

\begin{itemize}
  \item Finish 20-poem gold set.
  \item Create review rubric.
  \item Create review JSONL records.
  \item Install Ollama and Qwen3.
  \item Test one local prompt manually.
\end{itemize}

\subsection{Weeks 5--6}

\begin{itemize}
  \item Build Python schemas with Pydantic.
  \item Build JSONL loader.
  \item Build first draft generation script.
  \item Save raw model outputs.
  \item Review 10 machine drafts.
\end{itemize}

\subsection{Weeks 7--8}

\begin{itemize}
  \item Add BGE-M3 embeddings.
  \item Build retrieval over glossary and context.
  \item Add retrieved context to prompts.
  \item Compare no-retrieval vs retrieval outputs.
  \item Improve prompt template.
\end{itemize}

\subsection{Weeks 9--10}

\begin{itemize}
  \item Create website project.
  \item Build poem page template.
  \item Build poet page.
  \item Build glossary page.
  \item Export JSON data for website.
\end{itemize}

\subsection{Weeks 11--12}

\begin{itemize}
  \item Build timeline page.
  \item Add search.
  \item Add source and credits page.
  \item Add contribution page.
  \item Deploy to GitHub Pages.
\end{itemize}

\section{First MVP Definition of Done}

The MVP is done when:

\begin{itemize}
  \item one poet has at least 20 reviewed works;
  \item every work has original text, source, transliteration, translation, and
  tashreeh;
  \item at least 20 glossary terms exist;
  \item at least 10 timeline events exist;
  \item the local AI draft script works for one poem;
  \item generated drafts are saved and reviewable;
  \item the static website shows poet, works, poem pages, glossary, and sources;
  \item no private or copyrighted material is committed accidentally.
\end{itemize}

\section{Common Mistakes to Avoid}

\begin{itemize}
  \item Starting with model training before building the corpus.
  \item Scraping modern translations without permission.
  \item Treating one English word as the final meaning of a mystical term.
  \item Mixing literal translation and tashreeh into one output.
  \item Letting the model invent biographical context.
  \item Building a chatbot before the archive is useful.
  \item Committing large model files to Git.
  \item Publishing uncertain claims without confidence labels.
  \item Making the website depend on a paid backend from the start.
\end{itemize}

\section{Recommended Source Links}

\begin{itemize}
  \item IndicTrans2: \url{https://github.com/AI4Bharat/IndicTrans2}
  \item IndicTrans2 model card: \url{https://huggingface.co/ai4bharat/indictrans2-indic-en-1B}
  \item NLLB documentation: \url{https://huggingface.co/docs/transformers/en/model_doc/nllb}
  \item Qwen3 model card: \url{https://huggingface.co/Qwen/Qwen3-8B}
  \item Ollama Qwen3: \url{https://ollama.com/library/qwen3}
  \item BGE-M3 embeddings: \url{https://huggingface.co/BAAI/bge-m3}
  \item WebLLM: \url{https://github.com/mlc-ai/web-llm}
  \item PEFT LoRA documentation: \url{https://huggingface.co/docs/peft/en/developer_guides/lora}
  \item Axolotl: \url{https://docs.axolotl.ai/}
  \item Ajab Shahar reference: \url{https://ajabshahar.com/map}
\end{itemize}

\section{Final Strategic Recommendation}

Build Abshaar as a public archive with AI assistance, not as a model-only
project. The most sustainable version is:

\begin{enumerate}
  \item public data and documentation;
  \item transparent human review;
  \item local open-weight models;
  \item static website for free access;
  \item optional local AI for advanced users;
  \item fine-tuned models only after the corpus is strong.
\end{enumerate}

This preserves the poetry, keeps the project open, and makes the work accessible
to readers who do not already know the languages or traditions.

\end{document}
